\section{Discussion}
\label{sec:discussion}

\subsection{What the experiments show}
\label{sec:discussion-findings}

The scaling experiments reproduce the central AGJK phenomenon. Across four
graph families the quantum search time fits $T_q\sim\HT^{0.44\text{--}0.52}$,
close to the predicted square-root law, and the peak success probability is
high enough that the fitted times are meaningful. The cycle gives the largest
speedup because its classical hitting time is almost quadratic. The complete
graph and hypercube have smaller finite-size speedups because their classical
hitting times grow almost linearly over the simulated range.

The multiple-target experiments refine the picture. Independent targets are
easier: spreading targets across the graph lowers both the classical hitting
time and the quantum search time. Clustered targets behave like a single
target, because reaching the cluster does not mean finding a particular vertex.
The extended hitting time is close to the classical hitting time for these
configurations. The dramatic $\HT^+\gg\HT$ gap appears only in the periodic
adversarial construction of AGJK. This is an important negative result for a
naive reading of $\HT^+$: multiple targets alone do not create the penalty.
The penalty is a property of the arrangement and of the resulting spectral
structure.

The semantic experiment shows a distance-dependent difference between CTQW and
CRW. For a near target the classical walk activates the neighbor immediately
and has a larger early AUC. For a far target the quantum walk maintains a
larger and later peak, and its activation is spread over a larger radius. This
is consistent with the intuition that coherent propagation can bridge
communities, while classical diffusion relaxes to the stationary distribution.
The result is a model-level prediction: a semantic-search experiment should
see a larger peak but a later response for cross-community targets under the
quantum model, and an earlier but lower activation under the classical model.

The decision experiment produces a context effect that the classical walk
cannot. Coupling two options shifts quantum choice toward the uncoupled
option, increasingly with the coupling strength, while the classical choice on
symmetric graphs stays exactly uniform. The effect is moderate (from $1/3$ to
about $0.5$ on the symmetric three-option graph). Under the detection readout
it is monotone and robust to laziness and edge-weight noise (spread
$\le0.045$); under random stopping it is noisier; under a first-threshold
readout it is swamped by the oscillation phase. This is the same lesson as in
E5: a coherent walk must be read out with a stopping rule that does not depend
on its phase.

The decoherence experiment answers the noise question only once the readout
is specified. Comparing rates at one deadline is unreliable, because the pure
walk oscillates and the comparison then measures the oscillation phase; our
own first analysis made this mistake and reported a $12\times$ gain that
disappears under deadline averaging. If measurement only resets the walk, it
makes the success curve flatter but not higher. If a measurement can also end
the search, the pure walk wins only at short deadlines that catch its first
peak; beyond that an intermediate rate $p\approx0.2$--$0.3$ beats both the pure
quantum and the classical walk for clustered targets, and for single targets
the best rate moves, non-monotonically, toward the classical limit. The
model therefore predicts that the useful amount of ``introspection'' depends
on how extended the target is and how much time is available, which is a
testable prediction for deadline-manipulation experiments.

\subsection{Relation to the literature}
\label{sec:discussion-literature}

The technical results are consistent with AGJK \citep{agjk2019}. Our numerical
protocol tests the simpler interpolated-walk algorithm (Conjecture~3) and cites
the fast-forwarding theorem as the proved guarantee. The adversarial torus
reproduces the qualitative $\HT^+\gg\HT$ behavior and shows that the effect
survives at modest sizes ($L=18$--$36$). We do not claim a new theorem; the
contribution is the numerical validation and the cognitive interpretation.

The semantic experiment is closest to Pavone and Faro \citep{pavone2025}, but
it uses a discrete marked-vertex search perspective rather than an insight
task. The decision experiment is related to Huang and Busemeyer
\citep{huang2022} and to the evidence-accumulation literature
\citep{chen2022}. The decoherence model is a simplified version of the
open-system ideas in Epping et al.\ \citep{epping2023}: projective
measurements at a fixed rate rather than a Lindblad master equation. The
qualitative conclusion---that noise can help at intermediate levels and hurt at
strong levels---is consistent with the broader quantum-cognition literature on
decoherence-assisted processes.

\subsection{Algorithmic versus biological claims}
\label{sec:discussion-claims}

The project is deliberately agnostic about biology. Quantum walks are used as a
mathematical process model, not as a claim that the brain is a quantum computer.
The algorithmic advantage is a statement about the number of walk steps needed
to reach a target under a particular graph and marked set. Whether a biological
system can realize such a walk is a separate question. We keep the two claims
separate throughout: the QIS results concern complexity and spectra, and the
cognitive results concern the qualitative patterns that a process model can
produce.

This separation also limits the interpretation of the context effects. The
decision experiment shows that a quantum walk can produce a similarity-like
shift in choice under option coupling. It does not show that human similarity
effects arise from quantum interference. The value of the model is that it
generates distinct predictions: a classical random walk and a quantum walk on
the same graph produce different choice distributions and different response
times, so an experiment could in principle distinguish them.

\subsection{Limitations}
\label{sec:discussion-limitations}

Several limitations should be stated clearly.
\begin{itemize}[leftmargin=1.4em]
  \item \textbf{No human data.} The experiments are numerical. The cognitive
  quantities are model-defined, not fitted to behavioral data.
  \item \textbf{Synthetic graphs.} The semantic graph is hand-built and small
  ($19$ vertices). It illustrates the mechanism but does not capture the scale
  or noise of real semantic memory.
  \item \textbf{Interpolation selection.} The parameter $s$ is chosen by a
  grid search on the AGJK lower bound. A different selection rule could change
  the constants in the search times, although the scaling exponents are stable.
  \item \textbf{Two walk formalisms.} The search experiments use discrete-time
  Szegedy walks, while the semantic experiment uses continuous-time walks. We
  compare them only within their own time normalizations and do not mix the
  two formalisms in a single fit.
  \item \textbf{Simplified noise.} The decoherence model uses projective
  position measurements and resets. It is not a full Lindblad or open-system
  model.
  \item \textbf{Finite sizes.} The graphs have at most a few hundred vertices.
  The asymptotic exponents are estimated from small ranges, especially for the
  grid and hypercube.
\end{itemize}

\subsection{Future work}
\label{sec:discussion-future}

Four extensions would strengthen the project. First, replace the toy semantic
graph with a ConceptNet subgraph and test whether the distance-dependent
activation pattern survives in a real semantic network. Second, fit the
decision model to behavioral choice and response-time data, using the coupling
parameter as a free similarity parameter. Third, replace the projective
measurement model with a Lindblad equation or a continuous dephasing model and
compare the optimal noise level. Fourth, implement the decision graph on a
small quantum device or circuit simulator to test whether the context effect
survives realistic noise. A larger graph study could also test how the size
of the context effect depends on branch lengths, coupling patterns, and the
measurement rate of the detection readout.
